Las organizaciones quieren saber a menudo qué obtendría un grupo de clientes si lo atendiera otro canal, con solo datos observacionales. La respuesta estándar repondera a las unidades del otro canal, con pesos de un clasificador calibrado. Mostramos cómo puede fallar esa respuesta cuando, como es típico de los datos de canal, cada canal registra las covariables a su manera. Una variable que solo registra un canal rompe el soporte por construcción; el clasificador ideal lo delata, y descartarla es seguro cuando no predice el resultado en el canal básico dadas las demás. Una variable medida con un instrumento distinto en cada canal es peor: el soporte puede sobrevivir, y el estimador converge a un valor desplazado que ningún análisis de los datos observados distingue de la verdad. Derivamos ese desplazamiento exactamente. Con datos públicos en los que existen los dos instrumentos para las mismas personas (índice de masa corporal declarado y medido en NHANES), el desplazamiento derivado coincide de forma aproximada con el sesgo observado, y los diagnósticos habituales de solapamiento mejoran en lugar de delatarlo. Por el camino, leemos el fallo de soporte en el clasificador mediante la descomposición de Lebesgue, mostramos que el estimador doblemente robusto extrapola en silencio donde falta soporte, damos pesos para grupos objetivo definidos por un corte en una puntuación conocida y adaptamos cotas de sensibilidad para la confusión no observada. Una referencia experimental (NSW/CPS) y una simulación cuyas predicciones se registraron antes de ejecutarla (informamos de las que fallaron) completan la evidencia.
